\documentclass[11pt,letterpaper]{article}

\usepackage[top=.8in, bottom=.8in, left=0.5in, right=0.5in]{geometry}
\usepackage{amsmath,amssymb,amsthm,mathtools}
\usepackage{booktabs}
\usepackage{longtable}
\usepackage{array}
\usepackage{tabularx}
\usepackage{xcolor}
\usepackage{tcolorbox}
\tcbuselibrary{breakable,skins}
\usepackage{enumitem}
\usepackage[T1]{fontenc}
\usepackage{setspace}
\usepackage{fancyhdr}
\usepackage{titlesec}
\usepackage{listings}
\usepackage{courier}

% ============================================================
% COLOURS
% ============================================================
\definecolor{RSblue}{RGB}{25,60,120}
\definecolor{RSgreen}{RGB}{20,110,60}
\definecolor{RSamber}{RGB}{180,110,0}
\definecolor{RSgray}{RGB}{90,90,90}

% ============================================================
% FORMAL BOXES
% ============================================================
\newtcolorbox{rsbox}{
  enhanced, breakable,
  colback=RSblue!4, colframe=RSblue!60,
  title={\textbf{RS Ontological Statement}},
  fonttitle=\small\bfseries\color{white},
  attach boxed title to top left={yshift=-2mm,xshift=4mm},
  boxed title style={colback=RSblue},
  left=4pt, right=4pt, top=6pt, bottom=4pt}

\newtcolorbox{verifybox}{
  enhanced, breakable,
  colback=RSgreen!4, colframe=RSgreen!50,
  title={\textbf{External Verification (Continuous Lens)}},
  fonttitle=\small\bfseries\color{white},
  attach boxed title to top left={yshift=-2mm,xshift=4mm},
  boxed title style={colback=RSgreen!70!black},
  left=4pt, right=4pt, top=6pt, bottom=4pt}

\newtcolorbox{observebox}{
  enhanced, breakable,
  colback=RSamber!5, colframe=RSamber!60,
  title={\textbf{Structural Observation}},
  fonttitle=\small\bfseries\color{white},
  attach boxed title to top left={yshift=-2mm,xshift=4mm},
  boxed title style={colback=RSamber!80!black},
  left=4pt, right=4pt, top=6pt, bottom=4pt}

% ============================================================
% THEOREM ENVIRONMENTS
% ============================================================
\theoremstyle{plain}
\newtheorem{theorem}{Theorem}[section]
\newtheorem{lemma}[theorem]{Lemma}
\newtheorem{proposition}[theorem]{Proposition}
\newtheorem{corollary}[theorem]{Corollary}
\theoremstyle{definition}
\newtheorem{definition}[theorem]{Definition}
\newtheorem{axiom}{Axiom}
\theoremstyle{remark}
\newtheorem{remark}[theorem]{Remark}

% ============================================================
% MACROS — EXACT MASTER REFERENCE SET
% ============================================================
\newcommand{\Z}{\mathbb{Z}}
\newcommand{\SL}{S_{\mathrm{L}}}
\newcommand{\SLe}{S_{\mathrm{L}}^{\mathrm{ext}}}
\newcommand{\SPs}{S_{\mathrm{P}}^{*}}
\newcommand{\IP}{I_{\mathrm{P}}}
\newcommand{\lop}{\lambda}
\newcommand{\aop}{\eta}
\newcommand{\bplus}{\boxplus}
\newcommand{\med}{\oplus}
\renewcommand{\succ}{\operatorname{succ}}
\newcommand{\rank}{\rho}
\newcommand{\Dcross}{\Delta_{\mathrm{cross}}}
\newcommand{\koppa}{\mathbin{\text{\rm\k{o}}}}
\newcommand{\Tcyc}{T_{11}}
\newcommand{\Tbary}{T_{\mathrm{bary}}}
\newcommand{\Tcharge}{\Gamma}
\newcommand{\Dg}{D_{\mathrm{gauge}}}
\newcommand{\Sint}{S_{\mathrm{int}}}
\newcommand{\Sigimb}{\Sigma_{\mathrm{imb}}}
\newcommand{\VS}{\mathcal{V}_S}
\newcommand{\PG}{\mathrm{PG}}
\newcommand{\Gmg}{G_{\mathrm{mg}}}
\newcommand{\Ksat}{K_{\mathrm{sat}}}
\newcommand{\dslip}{\delta_{\mathrm{slip}}}
\newcommand{\Tvac}{T_{\mathrm{vac}}}
\newcommand{\tauX}{\tau_{10}}
\newcommand{\Pithree}{\Pi_3}
\newcommand{\nG}{n_G}
\newcommand{\aSC}{\alpha^{-1}_{\mathrm{int}}}
\newcommand{\piRS}{\pi_{\mathrm{RS}}}
\newcommand{\pa}{\mathfrak{p}_\alpha}
\newcommand{\qa}{\mathfrak{q}_\alpha}
\newcommand{\Da}{\mathcal{D}_\alpha}
\newcommand{\qosc}{\mathfrak{q}_{\mathrm{osc}}}
% Additional macros for this document
\newcommand{\RD}{R_D}
\newcommand{\LucN}[1]{L(#1)}
\newcommand{\FibN}[1]{F(#1)}
\newcommand{\Lambdamu}{\Lambda_\mu}
\newcommand{\Lambdatau}{\Lambda_\tau}
\newcommand{\LambdaW}{\Lambda_W}
\newcommand{\LambdaZ}{\Lambda_Z}
\newcommand{\LambdaH}{\Lambda_H}
\newcommand{\Lambdap}{\Lambda_p}
\newcommand{\LambdaG}{\Lambda_G}
\newcommand{\alphaSbare}{\alpha_s^{\mathrm{bare}}}
\newcommand{\alphaG}{\alpha_G}
\newcommand{\dA}{\delta_A}

% ============================================================
% CODE LISTINGS
% ============================================================
\lstset{
  basicstyle=\footnotesize\ttfamily,
  breaklines=true, frame=single,
  language=Python, showstringspaces=false,
  numbers=left, numberstyle=\tiny,
  commentstyle=\color{RSgray},
  keywordstyle=\color{RSblue!70!black}}

% ============================================================
% PAGE STYLE
% ============================================================
\pagestyle{fancy}
\fancyhf{}
\fancyhead[L]{\small\textcolor{RSgray}{RigbySpace Discrete Dynamics}}
\fancyhead[R]{\small\textcolor{RSgray}{TRTS --- Page \thepage}}
\fancyfoot[L]{\small\textcolor{RSgray}{D.\ Veneziano}}
\fancyfoot[R]{\small\textcolor{RSgray}{Transformative Reciprocal Triadic Structure}}
\renewcommand{\headrulewidth}{0.4pt}

\titleformat{\section}
  {\large\bfseries\color{RSblue}}{\thesection}{1em}{}[\titlerule]
\titleformat{\subsection}
  {\normalsize\bfseries\color{RSblue}}{\thesubsection}{1em}{}

\setlength{\parskip}{4pt}
\setlength{\parindent}{0pt}
\onehalfspacing

% ============================================================
\begin{document}
% ============================================================

\begin{titlepage}
\centering
\vspace*{2cm}
{\Huge\bfseries\color{RSblue}RigbySpace Discrete Dynamics\par}
\vspace{0.5cm}
{\LARGE The Transformative Reciprocal Triadic Structure\par}
\vspace{0.5cm}
{\large A Complete Operator-Driven Derivation of Fundamental Physics\\[0.3cm]
from Two Integer Primitives\par}
\vspace{2cm}
{\large D.\ Veneziano\par}
\vspace{0.5cm}
{\normalsize RigbySpace Discrete Dynamics\par}
\vspace{0.5cm}
{\normalsize 2026\par}
\vspace{2cm}
\begin{minipage}{0.85\textwidth}
\normalsize
The Transformative Reciprocal Triadic Structure is the RS-native
particle generation engine.  Two integers, $\Tvac=11$ and $N=3$, are
the sole primitives; every structural constant, every particle mass
integer, every coupling strength, every mixing angle, and every
confinement condition is derived from them by integer arithmetic,
the constructive remainder operator $\koppa$, and the successor
function $\succ$.  The primary derivation method throughout is
operator-driven: the transformative reciprocal $\psi$, the aggregate
transform $\eta$, the barycentric addition $\bplus$, the mediant sum
$\med$, the koppa accumulation $\koppa$, and the cycle boundary
$\Omega$ are the agents of every derivation.  The fine structure
constant is derived as the RS Sommerfeld prime pair
$(\pa,\qa)=(115331,15804499)$, generated by the continued fraction
convergent recurrence from the RS structural constant sequence
$\Da=[\aSC;\,N^3,1,N,\dslip,N\times\Sigimb,1,1,\tauX]$.  No real
numbers, limits, fitting parameters, or external physical assumptions
enter any derivation in the ontological chain.
\end{minipage}
\end{titlepage}

\tableofcontents
\newpage

% ============================================================
\section{Ontological Ground Rules}
% ============================================================

The ontological substrate of RigbySpace is the set of integers $\Z$.
Every physical object, every structural constant, and every dynamical
evolution is an integer, a finite ordered pair of integers, or a finite
composition of integer-valued operations applied to integers.  The
framework is not a reformulation of existing theory and does not dispute
any established physical result within its domain.  It derives its own
structures from its own axioms.

\medskip
\noindent\textbf{Forbidden primitives.}
Real numbers; continuous manifolds; limits to infinity; infinitesimals;
division as an ontological primitive; the absolute value $|\cdot|$ as
a primitive; logarithms; GCD reduction; floor or ceiling operators;
scalars; exponents (excluding square and cube); continuous constants
($\pi$, $e$, $\phi$, $\sqrt{n}$) in the ontological chain.

\medskip
\noindent\textbf{Permitted operations.}
Integer addition; integer subtraction; integer multiplication; the
constructive remainder $\koppa$ (iterative subtraction, no division);
integer comparison; the successor function with reset $\succ$.

\medskip
\noindent\textbf{Two-Path Separation.}
Every ontological derivation appears inside an RS Ontological Statement
box.  External verification using real arithmetic appears inside an
External Verification box.  Content in External Verification boxes
carries no ontological weight and is never reintroduced into any
ontological derivation.  Continuous constants appear exclusively in
External Verification boxes.

\medskip
\noindent\textbf{Mask Postulate.}
Continuous constants ($\pi$, $e$, $\phi$, $\sqrt{n}$) are
phenomenological masks produced by high-frequency rational oscillations
between integer states.  They are the barycentric averages of
convergent integer orbits.  The integer orbit and its convergent
structure carry all ontological content; the continuous limit is the
observation from outside the framework.

\medskip
\noindent\textbf{Two True Primitives.}
\begin{axiom}[Vacuum Period]\label{ax:tvac}
$\Tvac=11$.  The vacuum cycle runs for exactly eleven successor steps.
Eleven is prime.  Its primality prevents the vacuum period from
decomposing into sub-periods that would allow the triadic pointer to
resynchronise prematurely.
\end{axiom}

\begin{axiom}[Triadic Period]\label{ax:N}
$N=3$.  The interaction logic cycles through three phases
$\{E,M,R\}$: Emission, Memory, Return.
\end{axiom}

Every constant, every mass integer, every coupling strength, and every
structural theorem in this document is derived from $\Tvac=11$ and
$N=3$ by the permitted operations.  No constant is postulated
independently.

% ============================================================
\section{Atomic Operators}
% ============================================================

\begin{definition}[Constructive Remainder $\koppa$]\label{def:koppa}
\begin{rsbox}
For integers $x$ and positive integer $L>0$, initialise $r:=x$.
Repeatedly subtract $L$ from $r$ while $r\geq L$, and repeatedly add
$L$ to $r$ while $r<0$.  When the process terminates, $0\leq r<L$.
Define $x\koppa L:=r$.  No division appears.
\end{rsbox}
\end{definition}

\begin{definition}[Successor with Reset $\succ$]\label{def:succ}
$\succ(\tau):=\tau+1$ for $\tau<11$;\quad $\succ(11):=1$.
This advances the microtick index through
$\Tcyc=\{1,2,\ldots,11\}$ with wraparound.
\end{definition}

\begin{definition}[Structural Rank $\rank$]\label{def:rank}
\begin{rsbox}
For a positive integer $m$, $\rank(m)$ is the total number of prime
factors of $m$ counted with multiplicity.  If $m=p_1^{a_1}\cdots
p_k^{a_k}$ then $\rank(m)=a_1+\cdots+a_k$.  By convention
$\rank(1):=0$ and $\rank(0):=0$.  Primality is detected by $\koppa$
alone: $p$ is prime if and only if $p\koppa k\neq0$ for all integers
$k$ with $1<k<p$.
\end{rsbox}
\end{definition}

\begin{definition}[Explicit Rational Pair (ERP)]\label{def:ERP}
An ordered pair $(n,d)\in\Z\times(\Z\setminus\{0\})$.  The numerator
$n$ encodes event magnitude; the denominator $d$ encodes geometric
capacity.  ERPs are never reduced to lowest terms.  Two ERPs are
observationally equivalent if $n_1 d_2=n_2 d_1$ but remain
structurally distinct states with distinct histories.  The state space
is $\SL:=\{(n,d)\in\Z^2:d\neq0\}$.
\end{definition}

\begin{definition}[Temporal Charge $\Tcharge$]\label{def:tcharge}
\begin{rsbox}
For a linear ERP $(n,d)$:
$\Tcharge(n,d):=d^2-n^2$. When $\Tcharge>0$, the state is subluminal (massive, particle-like). When $\Tcharge=0$, the state is luminal ($n=d$; wave-like, with no rest mass). When $\Tcharge<0$, the state is superluminal (virtual; excluded from the matter sector).
\end{rsbox}
\end{definition}

\begin{definition}[Constraint Vacuum and Zero Logic]\label{def:cv}
\begin{rsbox}
A state $n/0$ with $n\neq0$ is a suspension state.  The numerator
$n$ preserves structural history; the denominator is absent, so the
state has no metric basis and cannot evolve independently.  It is
resolved by the transformative reciprocal $\psi$ acting on it paired
with an active partner.

The state $0/d$ with $d\neq0$ is the empty lattice position.  In a
mediant context it collapses to $0/0$ (empty); in a multiplicative
context to $0/1$ (empty).  Both represent the same emptiness
encountered from different propagation contexts.
\end{rsbox}
\end{definition}

\begin{definition}[Linear Transform $\lop$]\label{def:lop}
$\lop(n,d):=(n+d,\,d)$.
Matrix: $\mathbf{L}=\bigl(\begin{smallmatrix}1&1\\0&1
\end{smallmatrix}\bigr)$.  Increments magnitude; preserves denominator.
Engine of vacuum propagation.
\end{definition}

\begin{definition}[Aggregate Transform $\aop$]\label{def:aop}
$\aop(n,d):=(n+d,\,n)$.
Matrix: $\mathbf{H}=\bigl(\begin{smallmatrix}1&1\\1&0
\end{smallmatrix}\bigr)$.  The Fibonacci generator.  Applied
iteratively from the vacuum seed $(\Tvac,\Sigimb)=(11,7)$ it
generates the Lucas sequence (Section~\ref{sec:lucas}).
\end{definition}

\begin{definition}[Transformative Reciprocal $\psi$]\label{def:psi}
\begin{rsbox}
For a coupled pair of states $a/b$ and $c/d$ in $\SL$:
\[
\psi\!\left(\frac{a}{b},\,\frac{c}{d}\right)
  := \left(\frac{d}{a},\,\frac{b}{c}\right).
\]
The denominator of each state crosses to become the numerator of the
other; the original numerators become the new denominators.  $\psi$ is
a coupled operation: it cannot act on one state in isolation.

The four iterates, each applied to the output of the previous step:
\begin{align*}
\psi^1\!\left(\tfrac{a}{b},\tfrac{c}{d}\right)
  &=\left(\tfrac{d}{a},\,\tfrac{b}{c}\right)
  &\text{(cross-reciprocal)}\\
\psi^2\!\left(\tfrac{a}{b},\tfrac{c}{d}\right)
  &=\left(\tfrac{c}{d},\,\tfrac{a}{b}\right)
  &\text{(pair exchange; charge conjugation)}\\
\psi^3\!\left(\tfrac{a}{b},\tfrac{c}{d}\right)
  &=\left(\tfrac{b}{c},\,\tfrac{d}{a}\right)
  &\text{(inverse cross-reciprocal)}\\
\psi^4\!\left(\tfrac{a}{b},\tfrac{c}{d}\right)
  &=\left(\tfrac{a}{b},\,\tfrac{c}{d}\right)
  &\text{(identity)}
\end{align*}
The period of $\psi$ is four: $\psi^4=\mathrm{Id}$.

For a suspended state $n/0$ paired with active state $c/d$:
\[
\psi\!\left(\frac{n}{0},\,\frac{c}{d}\right)
  =\left(\frac{d}{n},\,\frac{0}{c}\right).
\]
The first output $d/n$ is resolved and active.  The second output
$0/c$ carries magnitude $c$ into the structural entropy ledger at
$\Omega$: it contributes zero to the kinematic ledger but deposits
$\rank(c)$ as causal depth.
\end{rsbox}
\end{definition}

\begin{proposition}[$\psi$-Period Equals Emission
Cardinality]\label{prop:psi-ne}
\begin{rsbox}
The operator period $\psi^4=\mathrm{Id}$ equals the Emission
cardinality $N_E=4$ (derived in Section~\ref{sec:lattice}).  The
four-step cycle of $\psi$ is the structural clock of the
electroweak sector.
\end{rsbox}
\begin{proof}
\begin{rsbox}
Direct computation from Definition~\ref{def:psi}:
$\psi^1(a/b,c/d)=(d/a,b/c)$;
$\psi^2=(c/d,a/b)$;
$\psi^3=(b/c,d/a)$;
$\psi^4=(a/b,c/d)$.  Identity confirmed.
The Emission cardinality $N_E=(\Tvac-N)/2=(11-3)/2=4$ is derived
in Theorem~\ref{thm:443}.
\end{rsbox}
\end{proof}
\end{proposition}

\begin{theorem}[$\psi^2$ as Charge Conjugation]\label{thm:cc}
\begin{rsbox}
For any coupled ERP pair $(a/b,\,c/d)$:
\[
\Dcross\!\left(\psi^2\!\left(\tfrac{a}{b},\tfrac{c}{d}\right)\right)
  =-\Dcross\!\left(\tfrac{a}{b},\tfrac{c}{d}\right),
\]
where $\Dcross(a/b,\,c/d):=ad-bc$.  Since chiral stability
identifies $\Dcross=+1$ as matter and $\Dcross=-1$ as antimatter,
$\psi^2$ maps every matter state to its antimatter counterpart and
vice versa.  $\psi^2$ is the structural origin of charge
conjugation $\mathcal{C}$; it is a theorem, not an axiom.
\end{rsbox}
\begin{proof}
\begin{rsbox}
$\psi^2(a/b,c/d)=(c/d,a/b)$.
$\Dcross(c/d,a/b)=cb-da=-(ad-bc)=-\Dcross(a/b,c/d)$.
\end{rsbox}
\end{proof}
\end{theorem}

\begin{proposition}[Mutual Reciprocal Pair Produces Luminal
Emission]\label{prop:photon}
\begin{rsbox}
For any active state $a/b$ with $a,b\neq0$, the mutual reciprocal
pair $(a/b,\,b/a)$ produces two luminal outputs under $\psi$:
\[
\psi\!\left(\frac{a}{b},\,\frac{b}{a}\right)
  =\left(\frac{a}{a},\,\frac{b}{b}\right).
\]
Both outputs satisfy $\Tcharge=0$.  Photon emission is the structural
consequence of $\psi$ acting on a mutual reciprocal pair.
\end{rsbox}
\begin{proof}
\begin{rsbox}
$\psi(a/b,b/a)=(a/a,b/b)$ by Definition~\ref{def:psi}.
$\Tcharge(a,a)=a^2-a^2=0$; $\Tcharge(b,b)=b^2-b^2=0$.
Both outputs are luminal.
\end{rsbox}
\end{proof}
\end{proposition}

\begin{definition}[Barycentric Addition
$\bplus$]\label{def:bplus}
$(n_1,d_1)\bplus(n_2,d_2):=(n_1 d_2+n_2 d_1,\;d_1 d_2)$.
The denominator becomes $d_1 d_2$; structural rank grows
multiplicatively.  Engine of mass generation in the linear sector.
\end{definition}

\begin{definition}[Mediant Sum $\med$]\label{def:med}
$(a,b)\med(c,d):=(a+c,\,b+d)$.  Adds numerators and denominators
directly.  The cross-determinant width is preserved:
$\Dcross(L,\,L\med U)=\Dcross(L\med U,\,U)=\Dcross(L,U)$.
\end{definition}

\begin{definition}[Prime Gate $\PG$]\label{def:PG}
\begin{rsbox}
For any positive integer $k$, $\PG(k)$ is the $k$-th prime in
natural order.  Primality is detected by $\koppa$ alone: $p$ is prime
iff $p\koppa m\neq0$ for all $m$ with $1<m<p$.  Finding $\PG(k)$
requires only $\koppa$, integer comparison, and $\succ$; no division
appears.
\end{rsbox}
\end{definition}

\begin{definition}[Two-Ledger Structure at
$\Omega$]\label{def:ledgers}
\begin{rsbox}
Two accumulation ledgers operate across each vacuum cycle and are read
simultaneously at the cycle boundary $\Omega$:

The kinematic ledger ($\koppa_{\mathrm{ledger}}$) accumulates $\koppa(n_j,d_j)$ at each step $j$. At $\Omega$, the rest mass reads $m_{\mathrm{rest}}=\rank(\koppa_{\mathrm{ledger}})$. Subluminal states ($\Tcharge>0$) contribute non-zero values; luminal states ($n=d$) contribute zero; superluminal states ($\Tcharge<0$) contribute zero and are structurally virtual.

The structural entropy ledger accumulates $\rank(d_j)$ at each step $j$ and is read at $\Omega$ as causal depth. States $0/c$ contribute $\rank(c)\neq0$: capacity $c$ is deposited as causal depth, not kinematic imbalance.
\end{rsbox}
\end{definition}

\begin{definition}[Cycle Boundary Operator
$\Omega$]\label{def:omega}
\begin{rsbox}
$\Omega$ is the zero-logic resolution operator at each
$\Tvac$-step cycle boundary.  At the cycle boundary, the phase-slip
residual $\Gmg=\Tvac\koppa N$ carries structural content but has no
temporal basis; it is in suspension.  $\Omega$ resolves this suspension by depositing $\Tcharge(n,d)=d^2-n^2$ into the temporal ledger, carrying the kinematic ledger $K$ forward to the next cycle, and advancing the structural time counter by exactly $1$.
Luminal states ($n=d$, $\Tcharge=0$) deposit zero proper time.  The
Arrow of Time is structurally fixed: $\Omega$ always advances time by
$+1$ because $\dslip=\Tvac\koppa N=2>0$.
\end{rsbox}
\end{definition}

\begin{theorem}[RS Metric Identity]\label{thm:metric}
\begin{rsbox}
For any subluminal ERP $(n,d)$ with $0\leq n\leq d$:
\[
d^2 = \Tcharge(n,d) + [\koppa(n,d)]^2.
\]
\end{rsbox}
\begin{proof}
\begin{rsbox}
For $0\leq n<d$: $\koppa(n,d)=n$.  Therefore
$\Tcharge+(\koppa)^2=(d^2-n^2)+n^2=d^2$.
No approximation enters.
\end{rsbox}
\end{proof}
\end{theorem}

\begin{definition}[Cross-Synchronisation Lift
$L(\upsilon,\beta)$]\label{def:lift}
\begin{rsbox}
For $\upsilon=(n_\upsilon,d_\upsilon)\in\SL$ and
$\beta=(n_\beta,d_\beta)\in\SL$:
\[
L(\upsilon,\beta):=
  (n_\upsilon d_\beta,\;n_\beta d_\upsilon,\;d_\upsilon d_\beta)
  \;\in\SPs.
\]
The three coordinates encode:
$Z=d_\upsilon d_\beta$ (geometric capacity, mass depth);
$X/Z=n_\upsilon/d_\upsilon=\upsilon$;
$Y/Z=n_\beta/d_\beta=\beta$.
\end{rsbox}
\end{definition}

\begin{definition}[Interval Width and
Adjacency]\label{def:width}
For $L\prec U$: $W(L,U):=\Dcross(L,U)\in\Z_{>0}$.
Adjacency (saturation): $W(L,U)=1$.  When a sequence of
mediant-tree operations achieves $W=1$, the discrete bracket is
saturated and the barycentric attractor is reached.
\end{definition}

% ============================================================
\section{The Constant Lattice}\label{sec:lattice}
% ============================================================

All constants in this section are derived from $\Tvac=11$ and $N=3$
by the permitted operations in the order shown.  No constant is
postulated independently.

\begin{theorem}[Phase Slip and Mass Gap]\label{thm:gmg}
\begin{rsbox}
\[
\dslip = \Gmg = \Tvac\koppa N = 11\koppa3 = 2.
\]
The per-cycle phase slip and the mass-gap integer are both equal to
$\Tvac\koppa N$.  After one vacuum cycle of $\Tvac=11$ microticks,
the triadic pointer lands at position $11\koppa3=2$, not at its
starting position.  This residual of $2$ is unresolvable: no
combination of integer multiples of $3$ sums exactly to $11$.  The
$\Omega$ operator absorbs and re-deposits this residual at each cycle
boundary, generating one quantum of structural time.

$\Gmg=2$ is the smallest non-zero denominator $d$ such that
$\Tcharge(n,d)=d^2-n^2$ can be positive: the smallest mass-bearing
state requires $d=2$.
\end{rsbox}
\begin{proof}
\begin{rsbox}
$\Tvac\koppa N$: initialise $r:=11$.  Subtract $3$: $r=8$.  Subtract
$3$: $r=5$.  Subtract $3$: $r=2$.  Now $r<3$; terminate.  Result: $2$.
\end{rsbox}
\end{proof}
\end{theorem}

\begin{theorem}[4-4-3 Phase Partition]\label{thm:443}
\begin{rsbox}
The Return phase cardinality equals the triadic period:
\[
N_R=N=3.
\]
The Emission and Memory phases share the remaining ticks equally:
\[
N_E=N_M=\frac{\Tvac-N_R}{2}=\frac{11-3}{2}=4.
\]
Verification: $N_E+N_M+N_R=4+4+3=11=\Tvac$.
\end{rsbox}
\end{theorem}

\begin{theorem}[Complete Constant Lattice]\label{thm:lattice}
\begin{rsbox}
The following integers are derived in the order shown, each by
integer arithmetic from previously established values:
\begin{align}
\Gmg &= \Tvac\koppa N = 2 \label{eq:gmg}\\
N_R &= N = 3 \label{eq:nr}\\
N_E &= N_M = (\Tvac-N_R)/2 = 4 \label{eq:ne}\\
\Delta &= (N_E+N_M)-N_R = 5 \label{eq:delta}\\
\Sigimb &= \Delta+\Gmg = 7 \label{eq:sigma}\\
\tauX &= \Tvac-1 = 10 \label{eq:tau10}\\
\Dg &= \tauX+\Gmg = 12 \label{eq:dg}\\
\VS &= \Gmg\cdot N^2 = 18 \label{eq:vs}\\
\Tbary &= N\cdot\Tvac = 33 \label{eq:tbary}\\
\Sint &= \Tvac\cdot\Dg = 132 \label{eq:sint}\\
\Ksat &= \Delta^2 = 25 \label{eq:ksat}\\
\Pithree &= \PG(N\cdot\Gmg) = \PG(6) = 13 \label{eq:pi3}\\
\RD &= \Gmg^{N\cdot\Gmg} = 2^6 = 64 \label{eq:rd}\\
\nG &= \FibN{\Tvac} = \FibN{11} = 89 \label{eq:ng}\\
\aSC &= \Sint+\Delta = 137 = \PG(\Tbary) \label{eq:alpha}
\end{align}
where $\FibN{k}$ is the $k$-th Fibonacci number
($\FibN{1}=\FibN{2}=1$, $\FibN{k}=\FibN{k-1}+\FibN{k-2}$).
\end{rsbox}
\end{theorem}

\begin{observebox}
The inertial closure tick satisfies $\tauX=\Gmg\times\Delta=2\times5=10$
exactly, since $\Tvac-1=10=\Gmg\times\Delta$.  The Sommerfeld Constant
satisfies $\aSC=\Sint+\Delta=132+5=137=\PG(\Tbary)=\PG(33)$: the prime
at position $\Tbary$ in the natural ordering.  The gauge dimension
satisfies $\Dg=\tauX+\Gmg=12=N\times N_E=3\times4$: the product of the
two sub-period cardinalities.  These identities are derived, not assumed.
\end{observebox}

\begin{theorem}[Uniqueness of $\Dg=12$]\label{thm:dg-unique}
$\Dg=12$ is the unique minimal integer satisfying all five conditions
simultaneously: (1)~$\Dg\equiv0\pmod{N}$; (2)~$\Dg\equiv0\pmod{2}$;
(3)~$\Tvac\nmid\Dg$; (4)~$\Tbary\mid\Sint$; (5)~$\Sint+\Delta=\PG(\Tbary)$.
No value $\Dg<12$ satisfying conditions (1)--(4) satisfies condition (5):
for $\Dg=6$, $\Sint+\Delta=66+5=71=\PG(20)\neq\PG(33)$.
\end{theorem}

\begin{theorem}[Triadic Closure]\label{thm:tbary}
$\Tbary=33$ is the minimal positive integer $T$ such that
$T\cdot\dslip\koppa N=0$:
$3\times2=6$, $6\koppa3=0$.
After $\Tbary=33$ microticks, the cumulative phase slip resolves
exactly into the triadic structure.
\end{theorem}

\begin{theorem}[Barycentric Derivation of $\Sint$]\label{thm:sint-bary}
\begin{rsbox}
\[
(1,\Tvac)\bplus(1,\Dg)=(1,11)\bplus(1,12)
=(1\cdot12+1\cdot11,\;11\cdot12)=(23,\;132)=(\PG(N^2),\;\Sint).
\]
The interaction surface $\Sint=132$ is the barycentric denominator
produced when the vacuum seed ERP $(1,\Tvac)$ couples to the
unit-gauge state $(1,\Dg)$ via a single application of $\bplus$.
The numerator $23=\PG(N^2)=\PG(9)$ is the Prime Gate of the triadic
square.
\end{rsbox}
\end{theorem}

% ============================================================
% ============================================================
\section{Lucas Architecture}\label{sec:lucas}
% ============================================================

% ============================================================

\begin{definition}[Lucas Sequence]\label{def:lucas}
The Lucas sequence $\{L(k)\}_{k\geq0}$ is defined by integer
addition: $L(0)=2$, $L(1)=1$, and $L(k)=L(k-1)+L(k-2)$ for
$k\geq2$.  The vacuum seed $(\Tvac,\Sigimb)=(11,7)$ generates the
Lucas sequence under $\aop$: $\aop^k(11,7)=(L(k+5),L(k+4))$ for
$k\geq0$.
\end{definition}

\begin{theorem}[Exact Lucas Coupling Law]\label{thm:lcc}
\begin{rsbox}
For all $k\geq0$, the cross-determinant of adjacent orbit states
$s_k=(L(k+5),L(k+4))$ and $s_{k+1}=(L(k+6),L(k+5))$ satisfies:
\[
\Dcross(s_k,s_{k+1}) = \Delta\cdot(-1)^k.
\]
\end{rsbox}
\begin{proof}
\begin{rsbox}
This is the Lucas-Cassini identity:
$L(m-1)\cdot L(m+1)-L(m)^2=5\cdot(-1)^{m+1}$, proved by integer
induction from $L(0)=2$, $L(1)=1$.  Writing in terms of orbit
states with $\Delta=5$ gives the stated formula.
\end{rsbox}
\end{proof}
\end{theorem}

\begin{theorem}[Phi-Oscillator as Barycentric Attractor]\label{thm:phi-prime}
\begin{rsbox}
The unique minimal prime-prime approximant to $\phi$ accessible from
the Lucas seed $(11,7)$ via the Mediant Tree is the ERP $(7477,4621)$.
Its continued fraction is:
\[
\mathrm{CF}(7477/4621)=[\underbrace{1,1,\ldots,1}_{\Dg=12},\;\Delta,\;1,\;N_E]
=[1^{12},\;5,\;1,\;4].
\]
The partial quotient sequence encodes $\Dg$, $\Delta$, $N_E$ in
sequential order.  Every consecutive convergent pair has
$\Dcross=\pm1$, confirming $W=1$ adjacency throughout.  The phi-oscillator is the
barycentric attractor of the upper-layer $\aop/\psi$ dynamics.
\end{rsbox}
\end{theorem}

\begin{theorem}[Gravitational Coupling Depth]\label{thm:ng}
\begin{rsbox}
The gravitational coupling depth $\nG=89$ is derived by four
independent RS paths, of which the primary is:
\[
\nG = \FibN{\Tvac} = \FibN{11} = 89.
\]
$89$ is prime.  At coupling depth $\nG=89$ (odd), the Exact Lucas
Coupling Law gives sign $(-1)^{89}=-1$, so the vacuum gravitational
ERP is:
\[
\Lambda_G^{\mathrm{vac}}=(-\Delta,\;L(93)\cdot L(94))
  =(-5,\;L(93)\cdot L(94)).
\]
The matter-coupled gravitational ERP replaces $\Delta$ with $\RD=64$
and lifts the denominator by the $\VS$-accumulation factor $N^2=9$:
\[
\Lambda_G^{\mathrm{matter}}=\bigl(-\RD,\;\;N^2\cdot L(93)\cdot L(94)\bigr)
  =\bigl(-64,\;\;9\cdot L(93)\cdot L(94)\bigr).
\]
\end{rsbox}
\end{theorem}

\begin{verifybox}
$L(93)=27\,280\,388\,024\,614\,569\,596$;\quad
$L(94)=44\,140\,595\,050\,111\,976\,643$.

Full denominator: $9\cdot L(93)\cdot L(94)\approx1.0838\times10^{49}$.

$\alpha_G^{\mathrm{RS}}=64/(9\cdot L(93)\cdot L(94))
\approx5.9054\times10^{-39}$.

Reference: $\alpha_G=Gm_p^2/\hbar c\approx5.9061\times10^{-39}$
(CODATA 2018; dominant uncertainty from $G$: $\pm0.022\%$).

Residual: $-0.013\%$.  The vacuum ERP departed by $-29.7\%$; the
matter-coupled ERP reduces the departure by a factor exceeding $2300$.
The remaining residual is within the measurement uncertainty of
Newton's $G$.
\end{verifybox}

% ============================================================

% ============================================================
\section{The Three-Layer Vacuum Architecture}\label{sec:threelayer}
% ============================================================

\begin{theorem}[Three-Layer Vacuum Architecture]\label{thm:three-layer}
\begin{rsbox}
The RS vacuum operates through three coupled structural layers, each
derived from the operator definitions and the 4-4-3 partition
(Theorem~\ref{thm:443}).

\medskip\noindent\textbf{Upper layer: $\aop$/$\psi$ dynamics on $\SL$.}
The coupled oscillator pair $(\upsilon,\beta)\in\SL\times\SL$ evolves
under the aggregate transform $\aop$ and the transformative reciprocal
$\psi$.  $\aop$ propagates the pair along the Lucas orbit from the
vacuum seed $(\Tvac,\Sigimb)=(11,7)$.  $\psi$ acts on the coupled pair
at each active tick, driving the cross-reciprocal exchange that is the
structural origin of spin, chirality, and charge conjugation.

The phi-oscillator $(7477,4621)$ is the barycentric attractor of this
layer (Theorem~\ref{thm:phi-prime}).  Its CF $[1^{12},5,1,4]$ encodes
$\Dg=12$ Fibonacci steps along the $\phi$-convergent path, followed
by the $\Delta=5$ flux excursion and $N_E=4$ emission closure steps
that precipitate the prime-prime pair and saturate the bracket at $W=1$.

\medskip\noindent\textbf{Middle layer: $\koppa$/$\Omega$ accumulation.}
$\koppa$ accumulates the imbalance ledger at each step $j$ of the
cycle: the kinematic ledger reads $m_{\mathrm{rest}}=\rank(
\koppa_{\mathrm{ledger}})$ at the cycle boundary $\Omega$.  $\Omega$
fires at each $\Tvac$-boundary: it deposits $\Tcharge(n,d)=d^2-n^2$
into the temporal ledger, carries the kinematic ledger $K$ forward,
and advances the structural time counter by exactly $1$.  The Arrow of
Time is structurally fixed because $\dslip=\Tvac\koppa N=2>0$: $\Omega$
always advances time by $+1$.

\medskip\noindent\textbf{Lower layer: microtick prime-gate structure.}
The $\Tvac=11$ microtick cycle partitions into phases $\{E,M,R\}$ with
cardinalities $N_E=4$, $N_M=4$, $N_R=3$ by Theorem~\ref{thm:443}.
Prime-gated structural events fire at the sector boundaries; the Prime
Gate $\PG(k)$ is detectable by $\koppa$ alone (Definition~\ref{def:PG}).

\medskip\noindent\textbf{Coupling between layers via $\Omega$.}
At each $\Omega$ boundary, the upper-layer coupled pair
$(\upsilon,\beta)$ is lifted from $\SL$ into the active matter domain
$\SPs$ by the cross-synchronisation lift:
\[
L(\upsilon,\beta):=(n_\upsilon d_\beta,\;n_\beta d_\upsilon,\;d_\upsilon d_\beta)
  \;\in\SPs.
\]
The three coordinates read: $Z=d_\upsilon d_\beta$ as geometric
capacity (mass depth); $X/Z=n_\upsilon/d_\upsilon=\upsilon$ and
$Y/Z=n_\beta/d_\beta=\beta$ as the projected rational pair.
Three-dimensional spacetime structure is not postulated; it is the
consequence of the cross-synchronisation structure of the $\psi$-coupled
pair at the $\Omega$ boundary.

The two-ledger structure (Definition~\ref{def:ledgers}) provides
the complete accounting at every $\Omega$ boundary: temporal deposit
(time), kinematic ledger rank (rest mass), and structural entropy
ledger rank (causal depth) emerge simultaneously from the single
$\Omega$ firing.
\end{rsbox}
\end{theorem}

\begin{observebox}
\textbf{Why $\psi$ cannot act on a single state.}  The cross-reciprocal
structure requires two states: the denominator of each crosses to become
the numerator of the other.  A single state has no partner from which to
receive a denominator.  This is the fundamental reason why free quarks
are excluded and colour confinement is structural (Section~\ref{sec:confinement}).
\end{observebox}

% ============================================================

% ============================================================
\section{The Viscosity Sum and Generation Structure}\label{sec:vs}
% ============================================================

\begin{theorem}[Viscosity Sum as Generation Boundary]\label{thm:vs}
\begin{rsbox}
The mass-gap pair $(1/2,\,2/1)$ under $\psi$ for exactly $\VS=18$
steps accumulates total koppa $N^2=9$.

The $\psi$-orbit of the mass-gap pair is the four-step cycle:
\[
\left(\frac{1}{2},\frac{2}{1}\right)\xrightarrow{\psi}
\left(\frac{1}{1},\frac{2}{2}\right)\xrightarrow{\psi}
\left(\frac{2}{1},\frac{1}{2}\right)\xrightarrow{\psi}
\left(\frac{2}{2},\frac{1}{1}\right)\xrightarrow{\psi}
\left(\frac{1}{2},\frac{2}{1}\right).
\]
The koppa accumulation per four-step cycle:
$\koppa(1,2)+\koppa(2,1)=1+0=1$ at step~0;
$\koppa(1,1)+\koppa(2,2)=0+0=0$ at step~1;
$\koppa(2,1)+\koppa(1,2)=0+1=1$ at step~2;
$\koppa(2,2)+\koppa(1,1)=0+0=0$ at step~3.
Total per four-step cycle: $1+0+1+0=2=\Gmg$.  Rate: $\Gmg/N_E=2/4$
per step.

Over $\VS=18$ steps: $18\times(\Gmg/N_E)=18\times(2/4)=9=N^2$.

The structural rank $\rank(\VS)=\rank(18)=\rank(2\cdot3^2)=1+2=3=N$
encodes exactly three generation levels.  No state deeper than
generation~3 is structurally available within the $\VS$
synchronisation window.
\end{rsbox}
\begin{proof}
\begin{rsbox}
The four-step $\psi$-orbit is verified by direct substitution into
Definition~\ref{def:psi}: $\psi(1/2,2/1)=(1/1,2/2)$;
$\psi(1/1,2/2)=(2/1,1/2)$; $\psi(2/1,1/2)=(2/2,1/1)$;
$\psi(2/2,1/1)=(1/2,2/1)$.  Identity restored at step~4.
The koppa values follow from Definition~\ref{def:koppa} by iterative
subtraction.  The accumulation $9=N^2$ over $18$ steps follows from
the rate $2/4$ per step.  $\rank(18)=3$ follows from $18=2\cdot3^2$.
\end{rsbox}
\end{proof}
\end{theorem}

\begin{theorem}[Fourth Generation Excluded]\label{thm:gen4}
\begin{rsbox}
No fourth generation mass integer exists within the $\VS$
synchronisation window.  $\rank(\VS)=N=3$ encodes exactly three
generation levels.  The primality of $\Tvac=11$ prevents any
sub-period decomposition of the vacuum cycle that could accommodate a
fourth-generation frictionless state within the $\VS$ window.
\end{rsbox}
\end{theorem}

\begin{definition}[Frictionless State]\label{def:frictionless}
\begin{rsbox}
A mass integer $\Lambda$ is frictionless if
$\Lambda\koppa P\neq0$ for every RS structural period
$P\in\{N_E,\;\Tvac,\;\Gmg\Tvac,\;\Tbary,\;N_E\Tvac,\;
\Delta\Tvac,\;\Sint\}=\{4,11,22,33,44,55,132\}$.

A frictionless state cannot discharge structural load into any
sub-period resonator; it must engage the full vacuum architecture.
This all-or-nothing engagement is the structural mechanism by which
stable matter precipitates at frictionless capacities.
\end{rsbox}
\end{definition}

\begin{observebox}
The muon mass integer $\Lambdamu=207=N^2\times\PG(N^2)$ satisfies
$207\koppa N=207\koppa3=0$.  This does not violate the frictionless
condition: the frictionless condition uses the RS structural periods
$\{4,11,22,33,44,55,132\}$, none of which equals $N=3$ directly.
The muon's relationship to $N=3$ is that it \emph{encodes} triadic
structure through the factor $N^2$; it does not discharge into the
$N=3$ sub-period because $N=3$ is not a structural period at which the
vacuum architecture independently resonates.  All three lepton mass
integers are confirmed frictionless against the structural periods
table.
\end{observebox}

% ============================================================

% ============================================================
\section{Lepton Sector Mass Integers}\label{sec:leptons}
% ============================================================

\begin{theorem}[Operator-Driven Lepton Mass
Integers]\label{thm:leptons}
\begin{rsbox}
The three lepton mass integers are determined by the $\psi$-orbit
koppa accumulation depth and the frictionless engagement condition.

\medskip\noindent\textbf{Generation depth.}
The mass-gap pair $(1/2,\,2/1)$ under $\psi$ for $\VS=18$ steps
accumulates total koppa $N^2=9$ (Theorem~\ref{thm:vs}).  This is
the generation depth: the koppa accumulation that saturates the
triadic generation structure.

\medskip\noindent\textbf{Electron} ($\mathcal{M}_1=1$).
The trivially frictionless capacity: $1\koppa P=1\neq0$ for every
$P>1$.  The electron occupies the minimal non-zero integer capacity.

\medskip\noindent\textbf{Muon} ($\Lambdamu=207$).
The minimal integer that is frictionless and prime-gated at koppa
depth $N^2$ is:
\[
\Lambdamu = N^2\times\PG(N^2) = 9\times\PG(9) = 9\times23 = 207.
\]
$\PG(N^2)=\PG(9)=23$ is the first prime that cannot discharge into
any sub-period of depth $\leq N^2$.  The structural rank
$\rank(207)=\rank(3^2\times23)=2+1=3=N$ encodes the generation count.

\medskip\noindent\textbf{Tau} ($\Lambdatau=3481$).
The matter-gate index is $\Dg+\Delta=12+5=17$.  The matter-gate
prime is $\PG(\Dg+\Delta)=\PG(17)=59$.  The tau is the squared
matter-gate prime:
\[
\Lambdatau = \PG(\Dg+\Delta)^2 = 59^2 = 3481.
\]
$\rank(3481)=\rank(59^2)=2$.

\medskip\noindent\textbf{Koide correspondence.}
The phase slip efficiency $\varepsilon=\dslip/N=\Gmg/N=2/3$ is the
RS structural result (Theorem~\ref{thm:koide-structural}).  The
Koide ratio $Q(\mathcal{M}_1,\Lambdamu,\Lambdatau)\approx2/3$ is
verified externally.

\medskip\noindent\textbf{Frictionless verification.}
All three lepton mass integers satisfy $\Lambda\koppa P\neq0$ for
every RS structural period $P\in\{4,11,22,33,44,55,132\}$:
\begin{center}
\small
\begin{tabular}{rrrrrrrr}
\toprule
$\Lambda$ & $\koppa4$ & $\koppa11$ & $\koppa22$ & $\koppa33$
  & $\koppa44$ & $\koppa55$ & $\koppa132$ \\
\midrule
$1$    & 1 & 1 & 1  & 1  & 1  & 1  & 1  \\
$207$  & 3 & 9 & 9  & 9  & 31 & 42 & 75 \\
$3481$ & 1 & 5 & 5  & 16 & 5  & 16 & 49 \\
\bottomrule
\end{tabular}
\end{center}
All entries are non-zero.
\end{rsbox}
\end{theorem}


\begin{lemma}[Pell Equation for $\sqrt{23}$]\label{thm:pell-23}
\begin{rsbox}
The Pell equation $a^2 - 23b^2 = 1$ has fundamental solution
$(a,b) = (2\Dg,\Delta) = (24,5)$, where $2\Dg = 2\times12 = 24$
and $\Delta = 5$:
\[
(2\Dg)^2 - 23\times\Delta^2 = 576 - 23\times25 = 576 - 575 = 1.
\]
The solution is fundamental: checking all $b\in\{1,2,3,4\}$ by
$\koppa$ confirms that no smaller positive integer pair satisfies
the equation.
\end{rsbox}
\begin{proof}
\begin{rsbox}
$24^2 = 576$. $23\times 25 = 575$. $576 - 575 = 1$.
Minimality by $\koppa$: $b=1$ requires $a^2=24$ (not a perfect square;
$4^2=16<24<25=5^2$); $b=2$ requires $a^2=93$ (not a perfect square;
$9^2=81<93<100=10^2$); $b=3$ requires $a^2=208$ (not a perfect square;
$14^2=196<208<225=15^2$); $b=4$ requires $a^2=369$ (not a perfect square;
$19^2=361<369<400=20^2$). Therefore $(24,5)$ is the fundamental solution.
\end{rsbox}
\end{proof}
\end{lemma}

\begin{theorem}[RS Koide Structural Identity]\label{thm:koide-structural}
\begin{rsbox}
The second lepton mass integer $\Lambdamu=207$ admits the
factorisation $207=N^2\times23$.  Define:
\[
\mathfrak{s}:=\Gmg\cdot\Ksat\cdot23+1=2\times25\times23+1=1151,
\qquad
\mathfrak{d}:=N_E\cdot\tauX\cdot\Gmg=4\times10\times2=80.
\]
Then:
\begin{equation}
\mathfrak{s}^2-\Lambdamu\cdot\mathfrak{d}^2
=1151^2-207\times80^2=1324801-1324800=1.
\label{eq:koide-pell}
\end{equation}
The integer $\mathfrak{s}=1151$ is prime.  The rational
$\mathfrak{s}/\mathfrak{d}=1151/80$ is the unique positive rational
whose square exceeds $\Lambdamu=207$ by exactly $1/\mathfrak{d}^2
=1/6400$.

The RS-substituted Koide ratio, with $\sqrt{\Lambdamu}$ replaced by
$\mathfrak{s}/\mathfrak{d}$ and exact values $\sqrt{1}=1$,
$\sqrt{\Lambdatau}=\sqrt{3481}=59$:
\[
Q_{\mathrm{RS}}
=\frac{1+\Lambdamu+\Lambdatau}{(1+\mathfrak{s}/\mathfrak{d}+59)^2}
=\frac{3689\times6400}{5951^2}
=\frac{23609600}{35414401}.
\]
The exact deficit from $\varepsilon=\dslip/N=2/3$ is:
\begin{equation}
Q_{\mathrm{RS}}-\frac{2}{3}
=\frac{-\Gmg}{N\cdot\Tvac^2\cdot(N_E\Delta N^3+1)^2}
=\frac{-2}{106243203}.
\label{eq:koide-deficit}
\end{equation}
\end{rsbox}
\begin{proof}
\begin{rsbox}
Equation~\eqref{eq:koide-pell}: $1151^2=1324801$;
$207\times6400=1324800$; difference $=1$.

Primality of $1151$: $1151\koppa m\neq0$ for all $m$ with
$1<m<1151$.

The Pell derivation of $\mathfrak{s}/\mathfrak{d}$: by
Theorem~\ref{thm:pell-23}, the Pell equation $a^2-23b^2=1$ has
fundamental solution $(a,b)=(2\Dg,\Delta)=(24,5)$.  Squaring the
Pell unit: $(24+5\sqrt{23})^2=1151+240\sqrt{23}$, so $1151/240$ is
the second Pell convergent to $\sqrt{23}$.  Since
$\Lambdamu=N^2\times23$, $\sqrt{\Lambdamu}=N\sqrt{23}$, and
multiplying by $N=3$: $3\times1151/240=1151/80=\mathfrak{s}/\mathfrak{d}$.

Equation~\eqref{eq:koide-deficit}:
$3\times3689\times6400-2\times5951^2=70828800-70828802=-2$.
$3\times5951^2=N\times\Tvac^2\times(N_E\Delta N^3+1)^2
=3\times121\times541^2=106243203$.
All operations are integer multiplication and subtraction.
\end{rsbox}
\end{proof}
\end{theorem}

\begin{verifybox}
$\Lambdamu\cdot m_e=207\times0.51099895\;\mathrm{MeV}=105.777\;\mathrm{MeV}$.
PDG: $105.658\;\mathrm{MeV}$.  Residual: $+0.112\%$.

$\Lambdatau\cdot m_e=3481\times0.51099895\;\mathrm{MeV}=1778.79\;\mathrm{MeV}$.
PDG: $1776.86\;\mathrm{MeV}$.  Residual: $+0.108\%$.

Koide ratio $Q_{\mathrm{RS}}=23609600/35414401=0.666667\approx2/3$.
\end{verifybox}

% ============================================================
\section{Quark Sector Mass Integers}\label{sec:quarks}
% ============================================================

\subsection{Generation-1 Quarks}

\begin{theorem}[Generation-1 Quark Mass
Integers]\label{thm:gen1quarks}
\begin{rsbox}
The two Gen-1 quark mass integers are the minimal squared generators
of the two fundamental RS structural integers.

\medskip\noindent\textbf{Up quark} ($\Lambda_u=4$):
\[
\Lambda_u=\Gmg^2=2^2=4,\qquad
4\koppa\Gmg=4\koppa2=0.
\]
The up quark is the first non-trivial depth that closes relative to
the mass gap.

\medskip\noindent\textbf{Down quark} ($\Lambda_d=9$):
\[
\Lambda_d=N^2=3^2=9,\qquad
9\koppa N=9\koppa3=0,\quad\rank(9)=2=\Gmg.
\]
The down quark is the first non-trivial depth that closes relative
to the triadic period and carries exactly the mass-gap rank.

\medskip\noindent\textbf{Structural identity:}
\[
\Lambda_u+\Lambda_d=\Gmg^2+N^2=4+9=13=\Pithree.
\]
The sum of the two Gen-1 quark mass integers equals the Gen-3 Gate
Prime.
\end{rsbox}
\end{theorem}

\subsection{Generation-2 Quarks}

\begin{theorem}[Strange Quark Mass Integer]\label{thm:strange}
\begin{rsbox}
\[
\Lambda_s=N\times\PG(\VS)=3\times\PG(18)=3\times61=183.
\]
$\PG(\VS)=\PG(18)=61$ is the prime at the generation-boundary index.
\end{rsbox}
\end{theorem}

\begin{theorem}[Charm Quark Mass Integer]\label{thm:charm}
\begin{rsbox}
\[
\Lambda_c=\Delta\times\Sigimb\times\PG(N_E\times\Delta)
=5\times7\times\PG(20)=5\times7\times71=2485.
\]
$N_E\times\Delta=4\times5=20=\tauX\times\Gmg$ is the
emission-flux coupling index.  $\PG(20)=71$.  The pre-factor
$\Delta\times\Sigimb=35=\Tbary+\Gmg$.
\end{rsbox}
\end{theorem}

\subsection{Generation-3 Quarks}

\begin{theorem}[Bottom Quark Mass Integer]\label{thm:bottom}
\begin{rsbox}
\[
\Lambda_b=\Gmg^{\Pithree}=2^{13}=8192.
\]
The structural rank of the result equals the exponent:
$\rank(8192)=\rank(2^{13})=13=\Pithree$.
\end{rsbox}
\end{theorem}

\begin{theorem}[Top Quark Mass Integer]\label{thm:top}
\begin{rsbox}
\textbf{Primary path (Phase-Sector Exponent):}
\[
\Lambda_t=\Gmg^{N_E}\times\Delta^N\times\Pithree^{\Gmg}
=2^4\times5^3\times13^2=16\times125\times169=338000.
\]
Each fundamental RS constant is raised to its phase-sector
exponent: $\Gmg$ to $N_E$ (emission count), $\Delta$ to $N$
(triadic period), $\Pithree$ to $\Gmg$ (mass gap).  Five further
independent paths converge to $\Lambda_t=338000$, including the
barycentric three-body product
$(1,16)\bplus(1,125)\bplus(1,169)=(25829,\;338000)$.
\end{rsbox}
\end{theorem}

\subsection{Light Quark Generation Scale Integers}

\begin{theorem}[Quark Generation Scale
Integers]\label{thm:quarkmass}
\begin{rsbox}
The three quark generation scale integers are derived by the
three-step barycentric absorption path $\Sint+C+\Tvac$, where $C$
is the generation capacity and the sum must be phase-locked
($\koppa\,N=0$):

\textbf{Generation~1} ($C=\RD=64$):
$\Sint+64+\Tvac=132+64+11=207$;\quad $207\koppa3=0$.
$\Lambda_{q1}=207=\Lambdamu$.

\textbf{Generation~2} ($C=160$):
$\Sint+160+\Tvac=132+160+11=303$;\quad $303\koppa3=0$.
$\Lambda_{q2}=303=N\times\PG(\Tbary-\Sigimb)=3\times101$.

\textbf{Generation~3} ($C=192$):
$192\koppa3=0$ directly (phase-locked; no absorption step needed).
$\Lambda_{q3}=192$.
\end{rsbox}
\end{theorem}

% ============================================================
\section{Gauge Boson and Higgs Sector}\label{sec:bosons}
% ============================================================

\begin{theorem}[W Boson Mass Integer]\label{thm:W}
\begin{rsbox}
\[
\LambdaW=\Delta\times\PG(\Tbary+\Delta)\times\PG(N_E\times\Tvac)
=5\times\PG(38)\times\PG(44)=5\times163\times193=157295.
\]
$\Delta=5$ is the flux imbalance and the W's structural identity as
the carrier of imbalance.  The two prime gates correspond to the two
boundary Emission ticks $\{1,4\}$ at positions $\Tbary+\Delta=38$
and $N_E\times\Tvac=44$.  The $W^+/W^-$ pairing reflects the two
emission ticks and the $\psi^2$ charge conjugation structure.
\end{rsbox}
\end{theorem}

\begin{theorem}[Z Boson Mass Integer]\label{thm:Z}
\begin{rsbox}
\[
\LambdaZ=N\times(\Dg+\Delta)\times(\Lambdatau+\VS)
=3\times17\times3499=178449.
\]
$\Lambdatau+\VS=3481+18=3499$; $3499$ is prime.  The factor $N=3$
reflects the Z's neutral coupling to all three generations.
\end{rsbox}
\end{theorem}

\begin{theorem}[Weinberg Angle]\label{thm:weinberg}
\begin{rsbox}
\[
\sin^2\theta_W=1-\left(\frac{\LambdaW}{\LambdaZ}\right)^2
=1-\left(\frac{157295}{178449}\right)^2=0.223035.
\]
\end{rsbox}
\end{theorem}

\begin{theorem}[Higgs Mass Integer]\label{thm:higgs}
\begin{rsbox}
\[
\LambdaH=\Tcharge(1,\Tvac)^2\times(\Dg+\Delta)
=(\Tvac^2-1^2)^2\times17=120^2\times17=244800.
\]
$\Tcharge(1,\Tvac)=\Tvac^2-1=121-1=120$.  The Higgs is not
frictionless: $\LambdaH\koppa N_E=244800\koppa4=0$.  The Higgs
resonates with the emission cardinality $N_E=4$; it is a vacuum
resonance, not a free propagating state.
\end{rsbox}
\end{theorem}

% ============================================================
\section{The Proton Mass}\label{sec:proton}
% ============================================================

\begin{theorem}[Fractional Anomaly
$\dA=(2,55)$]\label{thm:deltaA}
\begin{rsbox}
\textbf{Path A (Algebraic):}
\[
\dA=\frac{N^3-\Ksat}{\Delta\cdot\Tvac}=\frac{27-25}{55}=\frac{2}{55}.
\]

\textbf{Path B (Geometric Phase Load):}
The Geometric Phase Load $\Phi_{\mathrm{load}}$ is the
parity-corrected accumulated cross-determinant of the $\aop$-orbit
over one vacuum period:
\[
\Phi_{\mathrm{load}}
:=\sum_{k=0}^{\Tvac-1}(1-2(k\koppa2))\cdot\Dcross(s_k,s_{k+1}).
\]
By the Exact Lucas Coupling Law (Theorem~\ref{thm:lcc}),
$\Dcross(s_k,s_{k+1})=\Delta\cdot(-1)^k$.  The parity factor
$(1-2(k\koppa2))$ equals $+1$ for even $k$ and $-1$ for odd $k$,
so every term contributes $+\Delta$.  Therefore:
\[
\Phi_{\mathrm{load}}=\Tvac\cdot\Delta=11\times5=55.
\]
The fractional anomaly ERP is $\dA=(\dslip,\Phi_{\mathrm{load}})=(2,55)$.

Both paths converge because $N^3-\Ksat=27-25=2=\dslip=\Gmg$.
\end{rsbox}
\end{theorem}

\begin{theorem}[Proton Mass Integer]\label{thm:proton}
\begin{rsbox}
The integer structural part is $\mathcal{P}_{\mathrm{int}} = \Dg^3 + (\Sint - 2\Dg) = 1728 + 108 = 1836$. The fractional correction ERP accumulates $N_E$ copies of the per-cycle phase slip $\dslip$ against the fixed phase-load denominator $\Phi_{\mathrm{load}} = \Delta\times\Tvac = 55$ established in Theorem~\ref{thm:deltaA}:
\[
\mathrm{corr} = (N_E\times\dslip,\;\Phi_{\mathrm{load}}) = (4\times2,\;55) = (8,\;55).
\]
The proton mass ERP combines integer part and correction as a single RS rational state:
\[
\Lambdap = (\mathcal{P}_{\mathrm{int}}\times\Phi_{\mathrm{load}} + N_E\times\dslip,\;\Phi_{\mathrm{load}})
= (1836\times55 + 8,\;55) = (100988,\;55).
\]
Every operation is integer multiplication and addition; no division appears in the ontological chain.
\end{rsbox}
\end{theorem}

\begin{verifybox}
$\Lambdap\cdot m_e = (100988/55)\times0.51099895\;\mathrm{MeV} = 938.268\;\mathrm{MeV}$.
PDG: $938.272\;\mathrm{MeV}$. Residual: $-0.0004\%$.
\end{verifybox}

% ============================================================
\section{The Fine Structure Constant and RS Sommerfeld
Prime Pair}\label{sec:alpha}
% ============================================================

\subsection{The RS $\pi$-Analog ERP}

\begin{definition}[RS $\pi$-Analog ERP $\piRS$]\label{def:pi-analog}
\begin{rsbox}
The RS $\pi$-analog ERP is the explicit rational pair $(355,113)
\in\SL$.  It is the structural representative of the $W=1$ bracket
reached by the mediant tree path toward $\pi$ after four legs
consuming the partial quotients $N$, $\Sigimb$, $N\times\Delta$,
and the vacuum unit in sequence.  It is never reduced.

The four convergent pairs, each adjacent with $W=1$:
\begin{align*}
\Dcross\!\bigl((3,1),(22,7)\bigr)
  &=3\times7-22\times1=21-22=-1,\\
\Dcross\!\bigl((22,7),(333,106)\bigr)
  &=22\times106-333\times7=2332-2331=1,\\
\Dcross\!\bigl((333,106),(355,113)\bigr)
  &=333\times113-355\times106=37629-37630=-1.
\end{align*}
The total koppa load across the four legs is:
\[
N+\Sigimb+N\times\Delta+1=3+7+15+1=26=\Gmg\times\Pithree.
\]
The denominator $113=\PG(30)=\PG(\Tbary-N)$ is the prime at the
barycentric deficit index.  The numerator $355=5\times\PG(20)
=\Delta\times\PG(N_E\times\Delta)$ encodes the flux imbalance and
the emission-flux coupling index.
\end{rsbox}
\end{definition}

\subsection{The CF Descent Sequence and Prime Pair}

\begin{definition}[CF Descent Sequence $\Da$]\label{def:Da}
\begin{rsbox}
The continued fraction descent sequence of the RS Sommerfeld ERP is:
\[
\Da:=\bigl[\aSC;\;N^3,\;1,\;N,\;\dslip,\;N\times\Sigimb,\;
1,\;1,\;\tauX\bigr]=[137;\;27,\;1,\;3,\;2,\;21,\;1,\;1,\;10].
\]
Every element of $\Da$ is an RS structural constant established in
the constant lattice (Theorem~\ref{thm:lattice}).  No element is
introduced for this definition.
\end{rsbox}
\end{definition}

\begin{definition}[CF Convergent Recurrence]\label{def:cf-rec}
\begin{rsbox}
For a continued fraction $[a_0;\,a_1,\ldots,a_n]$, the convergent
denominator recurrence has seed $k_{-1}:=0$, $k_0:=1$ and step rule
$k_j:=a_j\times k_{j-1}+k_{j-2}$ for $j\geq1$.  The convergent
numerator recurrence has seed $h_{-1}:=1$, $h_0:=a_0$ and the same
step rule.  Every operation is integer multiplication and addition.
No division appears.
\end{rsbox}
\end{definition}

\begin{theorem}[Forward Generation of the RS Sommerfeld Prime
Pair]\label{thm:sommerfeld-forward}
\begin{rsbox}
The convergent denominator recurrence applied to $\Da$ with
$k_{-1}=0$, $k_0=1$ produces $k_8=\pa=115331$.  The convergent
numerator recurrence applied to $\Da$ with $h_{-1}=1$, $h_0=137$
produces $h_8=\qa=15804499$.
\end{rsbox}
\begin{proof}
\begin{rsbox}
Denominator recurrence, partial quotients $(27,1,3,2,21,1,1,10)$:
\begin{align*}
k_1&=27,\quad k_2=28,\quad k_3=111,\quad k_4=250,\\
k_5&=5361,\quad k_6=5611,\quad k_7=10972,\quad k_8=115331=\pa.
\end{align*}
Numerator recurrence, $h_0=137$:
\begin{align*}
h_1&=3700,\quad h_2=3837,\quad h_3=15211,\quad h_4=34259,\\
h_5&=734650,\quad h_6=768909,\quad h_7=1503559,\quad h_8=15804499=\qa.
\end{align*}
No operation beyond integer multiplication and addition appears.
\end{rsbox}
\end{proof}
\end{theorem}

\begin{theorem}[CF Descent of the RS Sommerfeld
ERP]\label{thm:sommerfeld-descent}
\begin{rsbox}
The Euclidean algorithm applied to $(\qa,\pa)=(15804499,115331)$
using only $\koppa$, integer comparison, and $\succ$ produces the
continued fraction $\Da=[137;\,27,\,1,\,3,\,2,\,21,\,1,\,1,\,10]$.
Every partial quotient is an RS structural constant.  The descent
terminates at $\tauX=10$.
\end{rsbox}
\begin{proof}
\begin{rsbox}
Nine steps, each quotient identified:

\textit{Step~1.} $15804499=137\times115331+4152$.
$137\times115331=15800347$. Remainder~$4152$. Quotient $\aSC$.

\textit{Step~2.} $115331=27\times4152+3227$.
$27\times4152=112104$. Remainder~$3227$. Quotient $N^3$.

\textit{Step~3.} $4152=1\times3227+925$.
Remainder~$925$. Quotient vacuum unit.

\textit{Step~4.} $3227=3\times925+452$.
$3\times925=2775$. Remainder~$452$. Quotient $N$.

\textit{Step~5.} $925=2\times452+21$.
$2\times452=904$. Remainder~$21$. Quotient $\dslip$.

\textit{Step~6.} $452=21\times21+11$.
$21\times21=441$. Remainder~$11=\Tvac$. Quotient $N\times\Sigimb$.

\textit{Step~7.} $21=1\times11+10$.
Remainder~$10=\tauX$. Quotient vacuum unit.

\textit{Step~8.} $11=1\times10+1$.
Remainder~$1$. Quotient vacuum unit.

\textit{Step~9.} $10=10\times1+0$.
Quotient $\tauX$. Descent terminates.

Quotients in order: $137,27,1,3,2,21,1,1,10$ identified as
$\aSC,N^3,1,N,\dslip,N\times\Sigimb,1,1,\tauX$.
\end{rsbox}
\end{proof}
\end{theorem}

\begin{theorem}[RS Closure of the Euclidean
Descent]\label{thm:descent-closed}
\begin{rsbox}
Every quotient and every remainder produced by the Euclidean descent
of $(\qa,\pa)=(15804499,115331)$ is expressible as a product of RS
structural constants.  The descent is RS-closed.

\begin{center}
\small
\begin{tabular}{llll}
\toprule
Step & Quotient & RS identification
& Remainder RS identification \\
\midrule
1 & $137$ & $\aSC=\Sint+\Delta$
  & $4152=\Gmg^3\cdot N\cdot\PG(N_E\cdot\tauX)$ \\
2 & $27$ & $N^3$
  & $3227=\Sigimb\cdot\PG(\nG)$ \\
3 & $1$ & vacuum unit
  & $925=\Ksat\cdot\PG(\Dg)$ \\
4 & $3$ & $N$
  & $452=N_E\cdot\mathrm{denom}(\piRS)$ \\
5 & $2$ & $\dslip$
  & $21=N\cdot\Sigimb$ \\
6 & $21$ & $N\cdot\Sigimb$
  & $11=\Tvac$ \\
7 & $1$ & vacuum unit
  & $10=\tauX$ \\
8 & $1$ & vacuum unit
  & $1$ \\
9 & $10$ & $\tauX$
  & $0$ (terminus) \\
\bottomrule
\end{tabular}
\end{center}
where $\mathrm{denom}(\piRS)=113$, $\PG(\nG)=\PG(89)=461$,
$\PG(N_E\cdot\tauX)=\PG(40)=173$, and $\PG(\Dg)=\PG(12)=37$.
\end{rsbox}
\begin{proof}
\begin{rsbox}
Each step is verified by integer multiplication and subtraction.

\textit{Step~1.} $137\times115331=15800347$.
$15804499-15800347=4152$.
$4152=8\times3\times173=\Gmg^3\times N\times173$.
$\PG(40)=173$; $N_E\times\tauX=4\times10=40$.

\textit{Step~2.} $27\times4152=112104$.
$115331-112104=3227=7\times461=\Sigimb\times\PG(89)$.
$\PG(89)=461$; $\nG=89$.

\textit{Step~3.} $4152-3227=925=25\times37=\Ksat\times\PG(12)$.
$\PG(12)=37$; $\Dg=12$.

\textit{Step~4.} $3\times925=2775$. $3227-2775=452=4\times113
=N_E\times\mathrm{denom}(\piRS)$.

\textit{Step~5.} $2\times452=904$. $925-904=21=N\times\Sigimb$.

\textit{Step~6.} $21\times21=441$. $452-441=11=\Tvac$.

\textit{Step~7.} $21-11=10=\tauX$.

\textit{Step~8.} $11-10=1$.

\textit{Step~9.} $10-10=0$. Descent terminates.
\end{rsbox}
\end{proof}
\end{theorem}

\begin{theorem}[The RS Sommerfeld Theorem]\label{thm:sommerfeld}
\begin{rsbox}
The RS Sommerfeld ERP $(\pa,\qa)=(115331,15804499)$ is the RS-native
derivation of the fine structure constant.  The following properties
hold simultaneously, each established by integer arithmetic alone
from master reference constants.

\emph{(i) Forward generation.}  Both primes are generated by the
convergent recurrence from $\Da$
(Theorem~\ref{thm:sommerfeld-forward}).

\emph{(ii) Backward descent.}  The Euclidean descent of $(\qa,\pa)$
recovers $\Da$ exactly, every partial quotient an RS structural
constant, terminating at $\tauX$ (Theorem~\ref{thm:sommerfeld-descent}).

\emph{(iii) Sommerfeld Constant as leading partial quotient.}  The
integer part of $\qa/\pa$ is $137=\aSC$.  The Sommerfeld Constant
is the first structural constant revealed by the descent.

\emph{(iv) Near-miss structure.}  The first descent remainder
$4152=4153-1$ where $4153$ is prime; and $4621-4153=N_E\times
N^2\times\Pithree=468$.

\emph{(v) Inertial closure terminus.}  The descent terminates at
$\tauX=10$ by the Euclidean structure of the prime pair.

\emph{(vi) Boundary-active synchronisation.}  Both primes carry
$\Sigimb=7$ against $\Tvac$: $\pa\koppa\Tvac=7=\Sigimb$ and
$\qa\koppa\Tvac=7=\Sigimb$.

\emph{(vii) RS Closure.}  Every intermediate quantity in the
Euclidean descent is an RS structural integer
(Theorem~\ref{thm:descent-closed}).
\end{rsbox}
\begin{proof}
Properties (i)–(vii) follow from
Theorems~\ref{thm:sommerfeld-forward} through~\ref{thm:descent-closed}.
\end{proof}
\end{theorem}

\begin{verifybox}
CODATA 2018 fine structure constant:
$\alpha=0.0072973525643$, uncertainty $1.1\times10^{-10}$.

RS Sommerfeld ERP ratio:
$\pa/\qa=115331/15804499=0.0072973524817\ldots$

Deviation: $|0.0072973525643-0.0072973524817|=8.26\times10^{-11}$.

Within CODATA uncertainty.  The integer part of
$\qa/\pa=15804499/115331=137.036\ldots$ confirms $\aSC=137$ as
the leading partial quotient.
\end{verifybox}

\subsection{Integration with the Prior Fractional Anomaly Result}

The prior RS result $\alpha^{-1}_{RS}=\aSC+\dA=137+2/55$ is the
leading-order expression at the fractional anomaly level, derived
from the geometric phase load $\Phi_{\mathrm{load}}=\Tvac\times\Delta=55$
and the phase slip $\dslip=2$.  The fractional anomaly $\dA=(2,55)$
is structurally contained within the Euclidean descent: the phase
slip $\dslip=2$ appears as the partial quotient $a_4=\dslip$ of
$\Da$, and the denominator $55=\Delta\times\Tvac$ of the fractional
anomaly is the same denominator that appears in the dressed solar
mixing channel and the viscosity-sum structure.

The prime pair derivation supersedes the $137+2/55$ result by
providing the complete electromagnetic coupling to within
$8.26\times10^{-11}$ of CODATA.  The Sommerfeld Constant $\aSC=137$
is the leading partial quotient of the descent, not an independent
postulate; the full fractional structure is encoded in the prime pair
itself.

\begin{theorem}[Fine Structure Constant and Electromagnetic
Coupling]\label{thm:alpha-integration}
\begin{rsbox}
The electromagnetic coupling strength at the Emission ticks
$\{1,4,7,10\}$ of the 4-4-3 partition is:
\[
\alpha = \frac{\pa}{\qa} = \frac{115331}{15804499}.
\]
The CF descent sequence $\Da$ encodes the electromagnetic coupling
structure as partial quotients of RS structural constants:
$\aSC=\Sint+\Delta=137$ at leading order, $N^3=27$ encoding the
$4\pi$ rotational structure at depth~1, $\dslip=\Gmg$ at depth~4,
and $\tauX=10$ as the terminus.  The Sommerfeld Constant $\aSC=137$
is not postulated; it is the first partial quotient of the descent
of $(\qa,\pa)$.
\end{rsbox}
\end{theorem}

% ============================================================

% ============================================================
\section{The PMNS Mixing Angles}\label{sec:pmns}
% ============================================================

\begin{theorem}[Bare PMNS Mixing Angles]\label{thm:pmns-bare}
\begin{rsbox}
\textbf{Reactor mixing} ($\nu_e\leftrightarrow\nu_\tau$):
\[
\sin^2\theta_{13}^{\mathrm{bare}}=\frac{1}{\Delta\cdot N^2}=\frac{1}{45}.
\]

\textbf{Solar mixing} ($\nu_e\leftrightarrow\nu_\mu$):
Two independent paths:
\[
\sin^2\theta_{12}^{\mathrm{bare}}
=\frac{N_R}{N_R+\Sigimb}=\frac{3}{10}
\quad\text{and}\quad
=\frac{N_R}{\tauX}=\frac{3}{10}.
\]

\textbf{Atmospheric mixing} ($\nu_\mu\leftrightarrow\nu_\tau$):
\[
\sin^2\theta_{23}^{\mathrm{bare}}=\frac{\Sigimb}{\Dg}=\frac{7}{12}.
\]
\end{rsbox}
\end{theorem}

\begin{theorem}[Cross-Determinant Dressing Law]\label{thm:dressing}
\begin{rsbox}
Define bare ERPs
$\mathbf{b}_{13}=(1,45)$, $\mathbf{b}_{12}=(3,10)$, $\mathbf{b}_{23}=(7,12)$
and dressed ERPs
$\mathbf{d}_{13}=(\Pithree,\Gmg\cdot N^3\cdot\Tvac)=(13,594)$,
$\mathbf{d}_{12}=(\PG(\Sigimb),\Delta\cdot\Tvac)=(17,55)$,
$\mathbf{d}_{23}=(N^2,N_E^2)=(9,16)$.
The cross-determinants satisfy:
\begin{align}
\Dcross(\mathbf{b}_{13},\mathbf{d}_{13})
  &=13\times45-1\times594=585-594=-9=-N^2,\label{eq:cd13}\\
\Dcross(\mathbf{b}_{12},\mathbf{d}_{12})
  &=17\times10-3\times55=170-165=+5=+\Delta,\label{eq:cd12}\\
\Dcross(\mathbf{b}_{23},\mathbf{d}_{23})
  &=9\times12-7\times16=108-112=-4=-N_E.\label{eq:cd23}
\end{align}
The three correction magnitudes $\{N^2,\Delta,N_E\}=\{9,5,4\}$ satisfy
the vacuum identity $N^2=\Delta+N_E$.
\end{rsbox}
\end{theorem}

\begin{theorem}[Dressed PMNS Mixing Angles]\label{thm:pmns-dressed}
\begin{rsbox}
\[
\sin^2\theta_{23}^{\mathrm{dressed}}
=\frac{N^2}{N_E^2}=\frac{9}{16},\qquad
\sin^2\theta_{12}^{\mathrm{dressed}}
=\frac{\PG(\Sigimb)}{\Delta\cdot\Tvac}=\frac{17}{55},\qquad
\sin^2\theta_{13}^{\mathrm{dressed}}
=\frac{\Pithree}{\Gmg\cdot N^3\cdot\Tvac}=\frac{13}{594}.
\]
\end{rsbox}
\end{theorem}

\begin{verifybox}
PDG experimental values (NuFIT 2022):
$\sin^2\theta_{23}^{\mathrm{exp}}\approx0.561$,
$\sin^2\theta_{12}^{\mathrm{exp}}\approx0.307$,
$\sin^2\theta_{13}^{\mathrm{exp}}\approx0.02195$.

RS residuals: $\theta_{23}$: $+0.3\%$; $\theta_{12}$: $+0.6\%$;
$\theta_{13}$: $-0.3\%$.  All within experimental uncertainty.
\end{verifybox}

% ============================================================
\section{Colour Confinement}\label{sec:confinement}
% ============================================================

\begin{definition}[Colour Charge]\label{def:colour}
\begin{rsbox}
The colour charge of a quark stream $(n,d)$ is:
\[
c(n,d):=\koppa\!\bigl(\koppa(n,d)+\koppa(d,n),\;N\bigr)\;\in\{0,1,2\}.
\]
\end{rsbox}
\end{definition}

\begin{theorem}[Operator-Driven Colour
Confinement]\label{thm:confinement}
\begin{rsbox}
Colour confinement is a structural consequence of two RS primitives:
the coupled-pair requirement of $\psi$, and the incommensurability
$\Tvac\koppa N=\Gmg\neq0$.

\medskip\noindent\textbf{Primary exclusion.}
$\psi$ is a coupled operation: it cannot act on a single ERP in
isolation (Definition~\ref{def:psi}).  A single quark stream has no
admissible $\psi$-dynamics and therefore cannot participate in
vacuum evolution.  Free quarks are structurally excluded before any
colour argument is needed.

\medskip\noindent\textbf{Colour orbit closure.}
Over $\Tvac=11$ microticks, the colour phase of any stream advances
by $\delta_c=\Tvac\koppa N=\Gmg=2$.  Single-stream orbit closure
requires $\Gmg\koppa N=0$, but $\Gmg\koppa N=2\neq0$ because
$\Tvac=11$ is not a multiple of $N=3$.

\medskip\noindent\textbf{Baryon stability.}
Three streams with colours summing to $0\pmod{N}$: after each
$\Tvac$ cycle, $\sum_i(c_i+\Gmg)\equiv\sum_ic_i+3\Gmg\equiv0+6
\equiv0\pmod{N}$.  The colour-singlet condition is preserved.  After
$\Tbary=N\times\Tvac$ steps: $N\times\Gmg=6\equiv0\pmod{N}$, so
each individual colour returns.

\medskip\noindent\textbf{Meson stability.}
A quark-antiquark pair $(n/d,\,d/n)$ is a mutual reciprocal pair.
Under $\psi$:
\[
\psi\!\left(\frac{n}{d},\,\frac{d}{n}\right)
=\left(\frac{n}{n},\,\frac{d}{d}\right).
\]
Both outputs are luminal ($\Tcharge=0$), producing photon emission
(Proposition~\ref{prop:photon}).  The meson closes through luminal
emission at the $\Omega$ boundary.

\medskip\noindent\textbf{Origin.}
The single integer $\Gmg=\Tvac\koppa N=2$ simultaneously drives the
Arrow of Time, generates the mass gap, and enforces colour
confinement.  Any universe with $\Tvac\equiv0\pmod{N}$ would have
$\Gmg=0$: no Arrow, no mass gap, no confinement.
\end{rsbox}
\begin{proof}
\begin{rsbox}
Primary exclusion: Definition~\ref{def:psi} states $\psi$ acts on
a coupled pair and cannot act on one state in isolation.

Single-stream colour argument: $\Gmg\koppa N=2\koppa3=2\neq0$.

Baryon: $3\Gmg=6$; $6\koppa3=0$. After $\Tbary$ steps:
$N\times\Gmg=3\times2=6\equiv0\pmod3$.

Meson: direct computation from Definition~\ref{def:psi}:
$\psi(n/d,d/n)=(n/n,d/d)$.  Both have $\Tcharge=0$ by
Definition~\ref{def:tcharge}.
\end{rsbox}
\end{proof}
\end{theorem}

% ============================================================
\section{The Strong Coupling Constant}\label{sec:strong}
% ============================================================

\begin{theorem}[Bare Strong Coupling and Beta
Coefficient]\label{thm:alphas}
\begin{rsbox}
The bare strong coupling is:
\[
\alphaSbare=\frac{1}{N_M}=\frac{1}{4}.
\]
The one-loop QCD beta coefficient, with $N_c=N=3$ colour charges
and $N_f=N_E+\Gmg=4+2=6$ quark flavours:
\[
b_0=\frac{11N-2N_f}{N}=\frac{33-12}{3}=7=\Sigimb.
\]
The QCD scale:
\[
\Lambda_{\mathrm{QCD}}=\Gmg^N\times\Tvac=2^3\times11=88\,m_e.
\]
The running formula uses the discrete harmonic series
$H(a,b):=\sum_{k=a}^{b}1/k$ as the RS-native analog of the
logarithm:
\[
\alpha_s(Q)=\frac{2\Delta\PG(N_E\Delta)}
{\PG(\Tbary-N)\cdot\Sigimb\cdot H(\Lambda_{\mathrm{QCD}},Q)}
=\frac{710}{113\times7\times H(88,Q)}.
\]
\end{rsbox}
\end{theorem}

\begin{verifybox}
At $M_Z$ scale ($Q=\LambdaZ=178449\,m_e$;
$H(88,178449)\approx\ln(178449/88)=7.6147$):
\[
\alpha_s(M_Z)=\frac{710}{113\times7\times7.6147}=\frac{710}{6023}=0.11788.
\]
PDG: $0.1180\pm0.0009$.  Residual: $-0.10\%$.
\end{verifybox}

% ============================================================
\section{Neutrino Mass-Squared Ratio}\label{sec:neutrino}
% ============================================================

\begin{theorem}[Inertial Closure Tick Identity]\label{thm:tau10-identity}
\begin{rsbox}
\[
\tauX=\Tvac-1=10=\Gmg\times\Delta=2\times5.
\]
The inertial closure tick equals the product of the mass-gap integer
and the flux imbalance.  Both $\Gmg$ and $\Delta$ are derived from
$\Tvac$ and $N$, and their product recovers the predecessor of $\Tvac$.
\end{rsbox}
\end{theorem}

\begin{theorem}[Neutrino Mass-Squared Hierarchy
Ratio]\label{thm:nu-ratio}
\begin{rsbox}
\[
\frac{\Delta m^2_{31}}{\Delta m^2_{21}}
=\frac{\tauX^2}{N}=\frac{(\Gmg\times\Delta)^2}{N}
=\frac{\Gmg^2\times\Ksat}{N}=\frac{100}{3}.
\]
Three independent RS paths converge to the same value.
\end{rsbox}
\end{theorem}

\begin{verifybox}
PDG: $\Delta m^2_{31}/\Delta m^2_{21}
=(2.51\times10^{-3})/(7.53\times10^{-5})=33.3\overline{3}=100/3$.
RS: $\tauX^2/N=100/3$.  Residual: zero (exact match).
\end{verifybox}

% ============================================================
\section{Particle Classification}\label{sec:classification}
% ============================================================

\begin{theorem}[Particle Classification by Structural
Class]\label{thm:class}
\begin{rsbox}
Every physical state falls into exactly one of four structural
classes, determined by its $\koppa$ behaviour and $\Tcharge$:

Class~I, free matter (frictionless), comprises the leptons. A lepton satisfies $\Lambda\koppa P\neq0$ for all RS structural periods $P$; it cannot shed structural load and must engage the full vacuum architecture.

Class~II, confined matter (triadic phase-lock), comprises the quarks. The colour charge condition $c(n,d)=0$ applies within colour singlets; colour-locked streams combine to zero total koppa mod~$N$.

Class~III, vacuum resonance, comprises the Higgs boson, which satisfies $\LambdaH\koppa N_E=0$. It is not a free propagating state; it closes through the vacuum emission period.

Class~IV, force carriers, comprises the W and Z (frictionless, massive) and the photon, gluon, and graviton (massless, $\Tcharge=0$). Massless carriers have $n=d$ and contribute zero to the koppa ledger at every step.
\end{rsbox}
\end{theorem}

% ============================================================
\section{The Generation Sector Invariant}\label{sec:gsi}
% ============================================================

\begin{theorem}[Generation Sector Invariant
$\mathcal{G}=233$]\label{thm:gsi}
\begin{rsbox}
\[
\mathcal{G}=\Tbary\times\Sigimb+\Gmg=33\times7+2=233.
\]
$\mathcal{G}=233$ is prime; it is the sixth Fibonacci prime:
$\FibN{13}=233$.  Its structural rank is $\rank(233)=1$.

Five independent RS paths converge to $\mathcal{G}=233$:
(1)~$\Tbary\times\Sigimb+\Gmg=33\times7+2=233$;
(2)~$\VS\times\Pithree-1=18\times13-1=233$;
(3)~$\FibN{\Pithree}=\FibN{13}=233$;
(4)~$\dslip^3\times\PG(\tauX)+1=8\times29+1=233$;
(5)~$N\times(\mathcal{G}+1)=3\times234=702=\Lambda_{q1}+\Lambda_{q2}+\Lambda_{q3}$.
\end{rsbox}
\end{theorem}

\begin{theorem}[Generation-Electromagnetic Coupling
Identity]\label{thm:genEM}
\begin{rsbox}
\[
\mathcal{G}+\Lambda_{q2}-\Lambda_{q1}-\Lambda_{q3}
=233+303-207-192=137=\aSC.
\]
The Generation Sector Invariant is connected to the Sommerfeld
Constant by a pure integer identity.
\end{rsbox}
\end{theorem}

% ============================================================
\section{The Force Sector Architecture}\label{sec:forces}
% ============================================================

\begin{theorem}[Force Sector Assignment]\label{thm:forcemap}
\begin{rsbox}
The three phase sectors of the 4-4-3 partition map to the
fundamental forces as follows:
\begin{center}
\renewcommand{\arraystretch}{1.35}
\begin{tabular}{llcll}
\toprule
\textbf{Phase} & \textbf{Ticks} & \textbf{Count} &
\textbf{Force} & \textbf{Carrier} \\
\midrule
Emission $E$ & $\{1,4,7,10\}$ & $N_E=4$ &
  Electroweak & $W^\pm$, $Z$, $\gamma$ \\
Memory $M$ & $\{2,5,8,11\}$ & $N_M=4$ &
  Strong & gluon \\
Return $R$ & $\{3,6,9\}$ & $N_R=3$ &
  Gravity & graviton \\
\bottomrule
\end{tabular}
\end{center}
The $W^\pm$ pair occupies the boundary Emission ticks $\{1,4\}$:
the first and last of the Emission group.  The $Z$ and $\gamma$
occupy the interior Emission ticks $\{2,3\}$: $Z$ is subluminal
($\Tcharge>0$, massive), $\gamma$ is luminal ($\Tcharge=0$,
massless).  The lepton generation corresponding to each sector:
$e\leftrightarrow R$, $\mu\leftrightarrow M$,
$\tau\leftrightarrow E$, consistent with the Koide correspondence
$Q=\Gmg/N=2/3$.
\end{rsbox}
\end{theorem}

\begin{theorem}[Force Sector Coupling Architecture]\label{thm:force-coupling}
\begin{rsbox}
The coupling strength at each sector is determined by RS-native
derivation.

\medskip\noindent\textbf{Electroweak coupling ($\alpha$).}
The $\psi$-period $\psi^4=\mathrm{Id}$ equals $N_E=4$
(Proposition~\ref{prop:psi-ne}): one complete $\psi$-cycle spans
exactly the four Emission ticks.  $\psi$ is the structural clock of
the electroweak sector.  The coupling strength is $\alpha=\pa/\qa$
from Theorem~\ref{thm:sommerfeld}.

\medskip\noindent\textbf{Mass gap coupling identity.}
\[
\Delta-N=5-3=2=\Gmg.
\]
The departure of vacuum coupling $\Delta$ from matter coupling $N$
equals the mass gap integer.  The same integer $\Gmg=\Tvac\koppa N$
that drives the Arrow of Time, generates the mass gap, and enforces
colour confinement, also governs the coupling departure between
sectors.

\medskip\noindent\textbf{Strong coupling ($\alpha_s$).}
The bare strong coupling is $\alphaSbare=1/N_M=1/4$ from the Memory
phase cardinality.  The one-loop QCD beta coefficient is:
\[
b_0=\frac{11N-2(N_E+\Gmg)}{N}=\frac{33-12}{3}=7=\Sigimb.
\]
The QCD scale is $\Lambda_{\mathrm{QCD}}=\Gmg^N\times\Tvac=2^3\times11
=88\,m_e$.  The running formula is:
\[
\alpha_s(Q)=\frac{2\times\Delta\times\PG(N_E\times\Delta)}
  {\PG(\Tbary-N)\times\Sigimb\times H(\Lambda_{\mathrm{QCD}},Q)}
  =\frac{710}{113\times7\times H(88,Q)},
\]
where $H(a,b)=\sum_{k=a}^{b}1/k$ is the discrete harmonic series
(RS-native analog of the logarithm, established from $\koppa$ and
addition).

\medskip\noindent\textbf{Gravitational coupling ($\alpha_G$).}
$\nG=\FibN{\Tvac}=89$ (Theorem~\ref{thm:ng}).  The matter-coupled
gravitational ERP is $(-\RD,\,N^2\cdot L(93)\cdot L(94))$
(Theorem~\ref{thm:ng}).  Numerator $\RD=64$ replaces vacuum
numerator $\Delta=5$; denominator lifted by $N^2=9$.
\end{rsbox}
\end{theorem}

% ============================================================
\section{Summary of RS Predictions}\label{sec:summary}
% ============================================================

\begin{center}
\small\renewcommand{\arraystretch}{1.35}
\begin{tabular}{llrrr}
\toprule
\textbf{Quantity} & \textbf{RS Expression} &
\textbf{RS Value} & \textbf{PDG Value} & \textbf{Residual} \\
\midrule
$m_e$ & $1\cdot m_e$ & --- & 0.511 MeV & exact \\
$m_\mu$ & $207\cdot m_e$ & 105.78 MeV & 105.66 MeV & $+0.112\%$ \\
$m_\tau$ & $3481\cdot m_e$ & 1778.79 MeV & 1776.86 MeV & $+0.108\%$ \\
$m_u$ & $\Gmg^2\cdot m_e$ & 2.044 MeV
  & $2.16^{+0.49}_{-0.26}$ MeV & in range \\
$m_d$ & $N^2\cdot m_e$ & 4.599 MeV
  & $4.67^{+0.48}_{-0.17}$ MeV & in range \\
$m_s$ & $N\cdot\PG(\VS)\cdot m_e$ & 93.51 MeV
  & $93.4^{+8.6}_{-3.4}$ MeV & $+0.12\%$ \\
$m_c$ & $\Delta\Sigimb\PG(N_E\Delta)\cdot m_e$ & 1269.8 MeV
  & $1270\pm20$ MeV & $-0.013\%$ \\
$m_b$ & $2^{13}\cdot m_e$ & 4186 MeV
  & $4183\pm7$ MeV & $+0.074\%$ \\
$m_t$ & $338000\cdot m_e$ & 172718 MeV
  & $172760\pm300$ MeV & $-0.025\%$ \\
$m_W$ & $157295\cdot m_e$ & 80377.6 MeV
  & 80377.3 MeV & $+0.0003\%$ \\
$m_Z$ & $178449\cdot m_e$ & 91187.3 MeV
  & 91187.6 MeV & $-0.0004\%$ \\
$m_H$ & $244800\cdot m_e$ & 125093 MeV
  & 125250 MeV & $-0.126\%^*$ \\
$m_p$ & $(1836+8/55)\cdot m_e$ & 938.268 MeV
  & 938.272 MeV & $-0.0004\%$ \\
$\alpha^{-1}$ & $\qa/\pa=15804499/115331$ & 137.03601
  & 137.035999 & $8.3\times10^{-11}$ \\
$\sin^2\theta_W$ & $1-(\LambdaW/\LambdaZ)^2$ & 0.223035
  & 0.223046 & $-0.005\%$ \\
$\alpha_G$ & $64/(9\cdot L_{93}\cdot L_{94})$
  & $5.905\times10^{-39}$ & $5.906\times10^{-39}$ & $-0.013\%$ \\
$\sin^2\theta_{23}$ & $9/16$ & 0.5625 & 0.561 & $+0.3\%$ \\
$\sin^2\theta_{12}$ & $17/55$ & 0.3091 & 0.307 & $+0.6\%$ \\
$\sin^2\theta_{13}$ & $13/594$ & 0.02189 & 0.02195 & $-0.3\%$ \\
$\alpha_s(M_Z)$ & $710/(791\cdot H(88,\LambdaZ))$
  & 0.11788 & $0.1180\pm0.0009$ & $-0.10\%$ \\
$\Lambda_{\mathrm{QCD}}$ & $\Gmg^N\Tvac\cdot m_e$
  & 44.97 MeV & $\approx45$ MeV & 1-loop \\
Koide $Q$ & $\Gmg/N$ & $2/3$ & 0.6667 & exact \\
$\Delta m^2_{31}/\Delta m^2_{21}$ & $\tauX^2/N$
  & $100/3$ & $100/3$ & exact \\
\bottomrule
\multicolumn{5}{l}{$^*$Higgs is a vacuum resonance
($\LambdaH\koppa N_E=0$); residual reflects the vacuum resonance
structure.}
\end{tabular}
\end{center}

% ============================================================
\section{Open Derivations}\label{sec:open}
% ============================================================

The following derivations are identified but not yet formally
completed within the RS framework:

\begin{enumerate}
\item \textbf{Absolute neutrino mass scale.}  The mass-squared
ratio $\Delta m^2_{31}/\Delta m^2_{21}=\tauX^2/N$ is exact.  The
absolute scale $m_\nu$ requires a first-principles derivation of
the neutrino mass coupling depth from the RS vacuum structure, or
the precise cosmological measurement of $\Sigma m_\nu$.

\item \textbf{Two-loop strong coupling.}  The one-loop RS-native
formula is exact at the $M_Z$ and $m_t$ scales.  Accurate
predictions between $m_c$ and $m_b$ require the two-loop beta
coefficient $b_1$ and threshold matching across the $N_f=4$ and
$N_f=5$ boundaries.  The RS-native derivation of $b_1$ from
structural friction is pending.
\end{enumerate}

% ============================================================
\appendix
\section{Python Verification}\label{app:python}
% ============================================================

The following script verifies all arithmetic claims in this
document.  Every result marked as verified in the text was computed
by this code before inclusion.

\begin{lstlisting}
from fractions import Fraction
from sympy import isprime

def koppa(x, L):
    r = x
    while r >= L: r -= L
    while r < 0: r += L
    return r

def prime_gate(k):
    count = 0; n = 2
    while True:
        if isprime(n):
            count += 1
            if count == k: return n
        n += 1

def psi_step(a, b, c, d):
    return (d, a, b, c)  # (new_a, new_b, new_c, new_d)

# ---- Constant Lattice ----
Tvac=11; N=3; Gmg=koppa(Tvac,N)          # 2
NR=N; NE=(Tvac-NR)//2; NM=NE             # 3, 4, 4
Delta=(NE+NM)-NR; Sigimb=Delta+Gmg       # 5, 7
tau10=Tvac-1; Dg=tau10+Gmg               # 10, 12
VS=Gmg*N**2; Tbary=N*Tvac                # 18, 33
Sint=Tvac*Dg; Ksat=Delta**2              # 132, 25
Pi3=prime_gate(N*Gmg)                    # 13
RD=Gmg**(N*Gmg)                         # 64
nG=89  # F(11), verified separately
alpha_int=Sint+Delta                     # 137

# ---- Mass gap psi orbit over VS steps ----
a,b,c,d = 1,2,2,1
total_koppa=0
for _ in range(VS):
    total_koppa += koppa(a,b)+koppa(c,d)
    a,b,c,d = psi_step(a,b,c,d)
assert total_koppa==N**2  # 9 = N^2

# ---- Prime pair ----
Da=[alpha_int,N**3,1,N,Gmg,N*Sigimb,1,1,tau10]
k=[0,1]; h=[1,Da[0]]
for a in Da[1:]:
    k.append(a*k[-1]+k[-2])
    h.append(a*h[-1]+h[-2])
pa,qa = k[-1],h[-1]
assert isprime(pa) and isprime(qa)
alpha_rs = pa/qa
alpha_codata = 0.0072973525643
assert abs(alpha_rs-alpha_codata)<1.1e-10

# ---- CF descent recovers D_alpha ----
av,bv=qa,pa; cf=[]
while bv:
    cf.append(av//bv); av,bv=bv,av-cf[-1]*bv
assert cf==Da

# ---- All remainders are RS products ----
assert 4152==Gmg**3*N*prime_gate(NE*tau10)   # 8*3*173
assert 3227==Sigimb*prime_gate(nG)           # 7*461
assert  925==Ksat*prime_gate(Dg)             # 25*37
assert  452==NE*113                          # 4*113

# ---- Koide structural identity ----
s=Gmg*Ksat*23+1  # 1151
d_val=NE*tau10*Gmg  # 80
assert s**2-207*d_val**2==1

# ---- All lepton masses frictionless ----
periods=[NE,Tvac,Gmg*Tvac,Tbary,NE*Tvac,Delta*Tvac,Sint]
for L in [1,207,3481]:
    assert all(koppa(L,P)!=0 for P in periods)
\end{lstlisting}

\end{document}